\section{The local periodic words}
\label{sec:local}

We now classify cyclic solutions of~\eqref{eq:signs}--\eqref{eq:offsets}
without reference to the size or parity of the center \(r\).

\begin{lemma}[Complete symbol classification]
\label{lem:symbol_classification}
Let \((\varepsilon_i,\delta_i)\) be a nonempty cyclic sequence with
\(\varepsilon_i\in\{-1,1\}\) and \(\delta_i\in\{-1,0,1\}\)
satisfying~\eqref{eq:signs}--\eqref{eq:offsets} for an integer \(s\).
For the coordinates \(u_i=\varepsilon_i r+\delta_i\), all possibilities
are repetitions and cyclic rotations of the following words:
\begin{equation}
 \begin{array}{c|l}
 s & \text{centered coordinate words}\\ \hline
 -1 &(1-r),\ (1+r)\\
 0 &(-r,-r,r,r)\\
 1 &(-r-1,r,r),\ (r-1,-r,-r)\\
 3 &(-r-1,r-1).
 \end{array}
 \label{eq:centered_words}
\end{equation}
There are no solutions for any other \(s\).
\end{lemma}
\begin{proof}
Equation~\eqref{eq:signs} has three consequences. At an offset
\(\delta_i=1\), both neighboring signs equal \(\varepsilon_i\).
At an offset \(\delta_i=-1\), both neighboring signs equal
\(-\varepsilon_i\). At an offset \(\delta_i=0\), the neighboring
signs are opposite. In particular, adjacent offsets 1 and \(-1\) are
impossible: the first would make the adjacent signs equal, while the
second would make them opposite. This argument also applies across
the cyclic boundary and when the period is one or two.

The range assertion in Proposition~\ref{prop:local_reduction}, which
also follows directly from~\eqref{eq:offsets}, leaves only six cases.

\paragraph{Case \(s=-2\).}
Equation~\eqref{eq:offsets} requires \(\delta_i^2+2\le2\) at every
index, so all offsets vanish. Substitution then gives \(0=2\), which
is impossible.

\paragraph{Case \(s=-1\).}
An offset \(-1\) would require both neighbors to have offset 1,
contradicting the forbidden adjacency. If a zero offset occurred, its
neighbors would sum to 1, so one of them would be 1. At that neighboring
offset, equation~\eqref{eq:offsets} would require both adjacent offsets
to be 1, contrary to the original zero. Thus all offsets are 1.
Equation~\eqref{eq:signs} makes all signs equal, giving the two
constant words \((1-r)\) and \((1+r)\).

\paragraph{Case \(s=0\).}
An offset \(-1\) would require neighboring offsets 0 and 1, forbidden
by adjacency. Among the remaining values 0 and 1, a zero offset has two
zero neighbors, whereas an offset 1 requires neighbors 0 and 1. An
offset 1 would therefore be adjacent to a zero that cannot neighbor it.
All offsets are zero, and the sign equation becomes
\(\varepsilon_{i+1}=-\varepsilon_{i-1}\). Its words repeat
\((u,v,-u,-v)\), where \(u,v\in\{-1,1\}\). The four choices of
\((u,v)\) are cyclic rotations of \((-1,-1,1,1)\), giving the single
centered word in~\eqref{eq:centered_words}.

\paragraph{Case \(s=1\).}
If an offset 1 occurred, its neighbors would sum to zero. Neither
could be \(-1\), so both would be zero. But a zero adjacent to 1 would
require its other neighbor to be \(-2\), impossible. The offsets lie
in \(\{-1,0\}\). Each \(-1\) has two zero neighbors, and each zero
has one neighbor \(-1\) and one zero. The offsets must therefore repeat
\((-1,0,0)\). Starting with sign \(u\) at the offset \(-1\), the sign
equation forces the following signs to be \(-u,-u,u\). The signs have
the same period three. The choices \(u=-1\) and \(u=1\) give,
respectively, \((-r-1,r,r)\) and \((r-1,-r,-r)\).

\paragraph{Case \(s=2\).}
An offset 1 would require neighbors \(-1\) and 0, again forbidden.
A zero offset would have two neighbors with offset \(-1\). The sign
rule at each neighbor forces its sign to be the negative of the center
sign, so the two neighboring signs agree. The sign rule at the center
zero requires them to be opposite. Thus neither 1 nor 0 can occur.
If all offsets were \(-1\), equation~\eqref{eq:offsets} would read
\(-2=-1\), also impossible.

\paragraph{Case \(s=3\).}
A zero offset would require its neighbors to sum to \(-3\), impossible.
Each remaining offset has two neighbors \(-1\), so offset 1 is
forbidden. All offsets equal \(-1\), and the signs alternate. This
gives \((-r-1,r-1)\).

These cases exhaust every cyclic solution, with no restriction on its
initially specified period.
\end{proof}

For integer centers, \(c=r^2+s=-A\), and the four allowed values of
\(s\) give Table~\ref{tab:even} for every \(A\le-13\). For the odd
branch, \(r=k+1/2\), \(x_i=u_i-1/2\), and
\begin{equation}
 A=-k(k+1)-1-s.
 \label{eq:odd_table_parameter}
\end{equation}
The same words give Table~\ref{tab:odd} for every \(A\le-17\).
To include small parameters in the existence assertion, one must verify
the words without using either large-parameter threshold.

\begin{lemma}[Existence and exact periods]
\label{lem:existence}
Every entry in Tables~\ref{tab:even} and~\ref{tab:odd} is an orbit for
its stated parameter and has the stated exact period. The within-row
coincidences are precisely those stated in the tables.
\end{lemma}
\begin{proof}
The centered recurrence can be checked for arbitrary real \(r\).
For \(c=r^2-1\), the constant coordinates satisfy
\[
 (1\pm r)^2-c=2(1\pm r).
\]
For \(c=r^2+3\), the two-cycle identities are
\[
 (-r-1)^2-c=2(r-1),\qquad (r-1)^2-c=2(-r-1).
\]
For \(c=r^2+1\), the first three-cycle satisfies
\[
 (-r-1)^2-c=2r,\qquad r^2-c=-1=(-r-1)+r,
\]
and the second satisfies
\[
 (r-1)^2-c=-2r,\qquad (-r)^2-c=-1=(r-1)+(-r).
\]
For \(c=r^2\), every squared coordinate in \((-r,-r,r,r)\) minus
\(c\) is zero, as is the sum of its two neighbors. These identities
prove existence for all listed integer values \(r=k\), or
half-integer values \(r=k+1/2\) followed by the shift \(-1/2\).

In the even branch the two-cycle entries differ by \(2k\), so they
are distinct for \(k\ge1\). Neither three-cycle word is constant at
any integer \(k\ge0\); their two orbits coincide when \(k=0\).
For \(k\ge1\), the first three-cycle has a repeated positive
coordinate \(k\), whereas the second has a repeated negative
coordinate \(-k\), so they are not cyclic rotations. The four-cycle
has neither period one nor period two for \(k\ge1\). The two fixed
words coincide only at \(k=0\).

In the odd branch the two-cycle entries differ by \(2k+1\), which is
nonzero for \(k\ge0\). The first three-cycle is always nonconstant.
The second is constant exactly at \(k=0\); its constant value is
\(-1\), which is the fixed point already listed at \(A=-2\).
For \(k\ge1\), the repeated coordinates in the two three-cycles are
\(k\) and \(-k-1\), so the orbits are distinct. The four-cycle has
neither period one nor two because \(-k-1\ne k\). Finally, \(-k\)
and \(k+1\) are always distinct fixed coordinates. A nonconstant
three-letter word has least period three, completing the checks.
\end{proof}
